\chapter{Technologies, Tools and Development Environment}\label{ch:tools}

Every version in this chapter was read from the tool itself in the environment that produced
the committed results. Where a version could not be verified, that is stated rather than
guessed.

\section{Target platform}

\begin{longtable}{@{}p{0.23\textwidth}p{0.16\textwidth}p{0.55\textwidth}@{}}
\toprule
\textbf{Item} & \textbf{Version} & \textbf{Role and why} \\
\midrule
\endhead
ESP32-WROOM-32 & DevKit v1 & The artefact under study. Chosen because it is the platform the
research question is about: dual-core, with the radio on the same die. \\
Xtensa LX6 cores & 2 cores & Core affinity is the only isolation primitive available; the
whole architecture depends on there being two. \\
arduino-esp32 & 2.0.17 and 3.3.12 & Both evaluated, because the buzzer driver API changed
between them and the firmware must build under each. Resource cost is reported for both
(Table~\ref{tab:rep-build}). \\
FreeRTOS & as shipped with the above & Supplies \code{xTaskCreatePinnedToCore}, queues and
semaphores. No version independent of the framework is pinned. \\
\bottomrule
\end{longtable}

\subsection{Firmware dependencies}

\begin{longtable}{@{}p{0.30\textwidth}p{0.62\textwidth}@{}}
\toprule
\textbf{Library} & \textbf{Purpose in this firmware} \\
\midrule
\endhead
\code{PubSubClient}            & MQTT 3.1.1 client: connection, subscription, publication,
Last Will \\
\code{WiFiClientSecure}        & TLS transport to the broker \\
\code{Adafruit\_SSD1306}, \code{Adafruit\_GFX} & Local OLED status display \\
\code{Preferences}             & Non-volatile storage for the armed state and the replay
counter \\
\code{mbedtls} (via the SDK)   & HMAC-SHA256 for command authentication \\
\bottomrule
\end{longtable}

\section{Verification and experimentation environment}

This is a separate environment from the target, and the distinction is material: every result
reported in Chapter~\ref{ch:results} except the broker experiment was produced here, on a host,
and not on the device.

\begin{longtable}{@{}p{0.21\textwidth}p{0.17\textwidth}p{0.56\textwidth}@{}}
\toprule
\textbf{Tool} & \textbf{Version} & \textbf{Role and where it appears} \\
\midrule
\endhead
GCC / g++ & 13.3.0 & Compiles the firmware logic against host stubs, and the harnesses in
\code{experiments/*.cpp}. Also supplies AddressSanitizer and UndefinedBehaviorSanitizer for
the fuzzing campaign. \\
Python & 3.13.18 & All analysis, figure generation and verification scripts. Statistics are
implemented with the standard library only, so no external package can alter a reported
value. \\
TLA+ / TLC & TLC2 2026.10.24.013619 (rev \code{2ef2af8}) & Exhaustive model checking of
\code{formal/AlarmFSM.tla} and \code{formal/Isolation.tla}, and of the deliberately defective
configuration that must fail. \\
Java & OpenJDK, as installed & Runtime for \code{tla2tools.jar} (4{,}502{,}259 bytes,
committed). \\
Mosquitto & 2.0.18 & The live broker for experiment E6. The only component in the evaluation
that is external software observed rather than simulated. \\
\code{mosquitto\_passwd} & with 2.0.18 & Generates the password file for per-client
authentication. \\
OpenSSL & as installed & Generates the private certificate authority and server certificate
for the broker testbed. \\
matplotlib & 3.11.2 & Figures 3 and 5. Configured through \code{experiments/figstyle.py} with
\code{pdf.fonttype~42}, so TrueType fonts are embedded rather than Type~3. \\
cairosvg & current at build & Converts the generated SVG diagrams to PDF and PNG. \\
Pillow & current at build & Measures text in the real font for the figure-fit check. \\
pdflatex & \TeX\ Live 2024 & Builds the manuscript and this report. \\
Pandoc & 3.1.3 & Produces the editable DOCX of this report. \\
Git & as installed & Version control; the commit hash on the title page identifies the exact
state. \\
GitHub Actions & --- & Continuous integration: re-runs the formal verification and a reduced
adversarial campaign on every push. \\
\bottomrule
\end{longtable}

\subsection{Why the firmware is tested against host stubs}

The firmware logic is compiled for the host against a set of stub headers in
\code{tests/host/stubs/}, which provide the Arduino, FreeRTOS, \code{WiFi}, \code{PubSubClient},
\code{Preferences}, \code{Wire} and display interfaces.

The reason is throughput. A campaign of 9,600 adversarial trials and 3,000,000 fuzzing
iterations is not feasible on the device at any useful speed, and the sanitizers that make the
fuzzing meaningful do not exist for the Xtensa target. Compiling for the host makes both
possible.

The cost is equally clear and is stated wherever a result from this environment is reported:
\textbf{the host harness advances a logical clock}. There is no radio, no true parallel
execution and no contention for shared silicon. Latency variance in this environment is
identically zero, which is why no timing conclusion is drawn from it anywhere in this report.

The firmware logic itself is \emph{not} modified for testing. The same \code{smart\_home.ino} is
compiled; only the headers it includes differ.

\section{Resource cost on the target}

Measured at build time by the Arduino toolchain, for the two framework versions. This is the one
measurement in the project taken by real tooling on the real target, and it describes the binary
rather than the running device.

\begin{table}[htbp]
\centering
\caption[Firmware resource cost]{Firmware resource cost, measured at build time.}
\label{tab:rep-build}
\begin{tabular}{@{}lrrrr@{}}
\toprule
Framework version & Flash (B) & of budget & RAM (B) & of budget \\
\midrule
arduino-esp32 2.0.17 & 793{,}433 & 60.5\% & 46{,}372 & 14.2\% \\
arduino-esp32 3.3.12 & 968{,}221 & 73.9\% & 49{,}876 & 15.2\% \\
\bottomrule
\end{tabular}
\end{table}

The increase under the newer framework is substantial — roughly 175\,kB of flash — and is a
property of the framework rather than of this firmware. It is reported because a reader
reproducing the build will see one figure or the other depending on which version they install.

\section{Simulation and modelling tools}

\begin{longtable}{@{}p{0.31\textwidth}p{0.61\textwidth}@{}}
\toprule
\textbf{Tool} & \textbf{Role} \\
\midrule
\endhead
Wokwi (\code{wokwi/diagram.json}) & Circuit definition for browser-based simulation of the
peripherals. Used during development for wiring verification; \textbf{no reported result comes
from it}. \\
\code{experiments/ablation.cpp} & The twelve-configuration harness. Models broker admission and
access control in \code{deliverGuarded}/\code{admit}. \\
\code{experiments/campaign.cpp} & The adversarial campaign harness, E1 and E2. \\
\code{experiments/fuzz\_cmd.cpp} & The command-parser fuzzer, run under both sanitizers. \\
\code{experiments/sensitivity.cpp} & The detector parameter sweep. \\
\bottomrule
\end{longtable}

\section{Code quality and verification infrastructure}

A deliberate choice was made to treat the verification infrastructure as part of the artefact
rather than as scaffolding, because four defects were found \emph{in the harnesses themselves}
during this work — two of which produced results that flattered the intended conclusion.

\begin{longtable}{@{}p{0.32\textwidth}p{0.60\textwidth}@{}}
\toprule
\textbf{Check} & \textbf{What it prevents} \\
\midrule
\endhead
\code{tests/host/run\_tests.sh} & 129 assertions over the firmware logic, including three
negative build-guard tests that require a defective configuration to fail to compile \\
\code{formal/run\_tlc.sh} & Three model-checking configurations; the third \emph{must} report a
violation, so a vacuous invariant cannot recur silently \\
\code{experiments/test\_conformance.py} & 19,529 firmware transitions checked against the
specification's transition relation, with a negative control \\
\code{experiments/test\_harness.py} & Metamorphic and differential properties of the harness:
determinism, seed sensitivity, monotone attack effort, mode separation, and agreement of the
HMAC payload format with an independent implementation \\
\code{experiments/corpus/test\_rule.py} & Eight controls for the static rule, including this
project's own pre-repair guard \\
\code{experiments/validate\_figures.py} & Recomputes every plotted value from the raw per-trial
records, not from the derived summary \\
\code{experiments/check\_figure\_type.py} & Fails the build if any figure label would print
below 6\,pt \\
\code{experiments/check\_figure\_fit.py} & Measures every label in the real font and fails on
clipping \\
\bottomrule
\end{longtable}

The last two exist because a defect was found that visual inspection cannot detect: every
diagram in an earlier revision printed its smallest label at roughly 4\,pt, which is a function
of the canvas width and the figure's scale factor rather than anything visible on screen.

\section{What the environment could not do}

Recorded here so that no reader infers a capability that did not exist.

\begin{itemize}[leftmargin=1.6em,itemsep=4pt]
  \item \textbf{No physical ESP32 hardware} was available in the environment that produced
        these results. There were no serial devices, no USB, no \code{arduino-cli} and no
        wireless interface. Appendix~\ref{app:hardware} records the probe that established
        this.
  \item \textbf{No RF-shielded enclosure}, which is a hard prerequisite for the radio-layer
        condition.
  \item \textbf{No licensed similarity-checking tool.}
  \item \textbf{No institutional access} to Scopus or Journal Citation Reports, which is why no
        journal quartile is asserted anywhere in this work.
\end{itemize}
